\section{Related work}
\label{sec:related}

We organize related work by methodological relationship to this study rather than
as a list of results.

\paragraph{Subword tokenization foundations.}
BPE for open-vocabulary neural machine translation \citep{sennrich2016bpe},
WordPiece \citep{schuster2012wordpiece}, subword regularization and Unigram
language-model tokenization \citep{kudo2018subword}, and SentencePiece
\citep{kudo2018sentencepiece} define the tokenizer families and the
open-vocabulary motivation; byte-level BPE \citep{radford2019gpt2} and
WordPiece-based encoders \citep{devlin2019bert} are widely deployed. We do not
introduce a new tokenizer algorithm; we vary the training corpus within a fixed
byte-level BPE design.

\paragraph{Byte-level BPE and implementation baselines.}
Byte-level BPE \citep{radford2019gpt2} and its production implementations
\citep{openai2026tiktoken,huggingface2026tokenizers} provide robust,
full-coverage tokenization. We use one such implementation for the controlled
tokenizers and treat production \texttt{tiktoken} encodings as observational
baselines.

\paragraph{Domain-specific tokenizers and vocabularies.}
Prior work builds domain-specific vocabularies or tokenizers for scientific
\citep{beltagy2019scibert}, legal \citep{chalkidis2020legalbert}, legal/financial
\citep{bommarito2025kl3m}, and code \citep{feng2020codebert} corpora, and studies
tokenizer choice for domains \citep{dagan2024tokenizer}. Notably, some
domain-adapted models instead reuse a general-domain tokenizer
\citep{lee2019biobert}, which is precisely the vocabulary-mismatch setting this
line of work motivates. This study targets cloud/technical documentation and
isolates \emph{compression} rather than downstream model quality.

\paragraph{Vocabulary adaptation for existing models.}
A separate line adapts or augments an existing model's tokenizer---adaptive
tokenization \citep{sachidananda2021adaptive}, adaptive BPE and vocabulary
adaptation in fine-tuning \citep{balde2024adaptivebpe,herold2025vocabulary},
general-to-specialized adaptation \citep{liyanage2026adaptbpe}, and
parameter-efficient vocabulary adaptation \citep{balde2026learning}. These methods
modify an existing vocabulary and often evaluate downstream tasks; we instead
train tokenizers from scratch on byte-matched corpora and report intrinsic
compression only, so we make no model-substitution claim.

\paragraph{Pre-tokenization and scaling choices.}
Pre-tokenization \citep{kudo2018sentencepiece,dagan2024tokenizer} and the amount of
tokenizer training data \citep{reddy2025diminishing} both affect the learned
inventory. We fix byte-level BPE design choices and vary corpus composition and
vocabulary size, observing diminishing sensitivity across 8K/16K/32K.

\paragraph{Fertility, compression, and fairness metrics.}
Token counts are corpus- and language-sensitive and can create inequities
\citep{petrov2023unfairness}, motivating metrics beyond raw fertility
\citep{nayeem2025strr} and studies of tokenizer choice
\citep{ali2024tokenizerchoice,rust2021tokenizer}. We adopt tokens per 1{,}000
UTF-8 bytes as the outcome and explicitly avoid downstream quality or fairness
claims.

\paragraph{Leakage and evaluation controls.}
Deduplication improves the integrity of train/eval separation for language
modeling \citep{lee2022dedup}. We apply analogous exact and near-duplicate
controls to tokenizer-compression evaluation.

\medskip
Taken together, prior work establishes that tokenizer inventories are
corpus-sensitive and can affect efficiency and downstream systems. Our
contribution is a controlled, byte-matched, leakage-controlled case study
measuring how that sensitivity manifests as \emph{compression} on technical
documentation, including the tradeoff between full specialization and mixed-domain
training. This work sits primarily in the ``trained from scratch on domain data''
and ``fertility/compression measurement'' categories above.
